\documentclass[../../../include/open-logic-section]{subfiles}

\begin{document}
\olfileid{sth}{card-arithmetic}{expotough}

\olsection[بعض التبسيطات]{بعض التبسيطات في أس الكاردينالات}

مع أن تعريف $\canonord$ كان معقداً بعض الشيء، فإن ثمرته نتيجة مفيدة عن
جمع الكاردينالات وضربها، هي \olref[simp]{cardplustimesmax}. غير أن أس
الكاردينالات ما فوق المتناهي لا يقبل تبسيطاً بمثل هذه المباشرة. ولبيان
السبب نبدأ بنتيجة توسع نمطاً مألوفاً من الحالة المتناهية (وإن كان
برهانها على درجة عالية من التجريد):

\begin{prop}\ollabel{simplecardexpo}
$\cardexpo{\cardfont{a}}{\cardfont{b} \cardplus \cardfont{c}} =
\cardexpo{\cardfont{a}}{\cardfont{b}} \cardtimes
\cardexpo{\cardfont{a}}{\cardfont{c}}$، و
$\cardexpo{(\cardexpo{\cardfont{a}}{\cardfont{b}})}{\cardfont{c}} =
\cardexpo{\cardfont{a}}{\cardfont{b} \cardtimes \cardfont{c}}$، لأي
كاردينالات $\cardfont{a}, \cardfont{b}, \cardfont{c}$.
\end{prop}

\begin{proof}
لإثبات الادعاء الأول، انظر في دالة
$f \colon (\cardfont{b}\disjointsum\cardfont{c}) \to \cardfont{a}$.
والآن «اقسمها» بتعريف $f_\cardfont{b}(\beta) = f(\beta, 0)$ لكل
$\beta \in \cardfont{b}$، و$f_\cardfont{c}(\gamma) = f(\gamma, 1)$
لكل $\gamma \in \cardfont{c}$. تكون الخريطة
$f \mapsto \tuple{f_{\cardfont{b}}, f_{\cardfont{c}}}$
!!a{bijection} من
$\funfromto{\cardfont{b} \disjointsum \cardfont{c}}{\cardfont{a}}$ إلى
$(\funfromto{\cardfont{b}}{\cardfont{a}} \times
\funfromto{\cardfont{c}}{\cardfont{a}})$.

ولإثبات الادعاء الثاني، انظر في دالة
$f \colon \cardfont{c} \to
(\funfromto{\cardfont{b}}{\cardfont{a}})$؛ فلكل
$\gamma \in \cardfont{c}$ دالة
$f(\gamma) \colon \cardfont{b} \to \cardfont{a}$. وعرّف الآن
$f^*(\beta, \gamma) = (f(\gamma))(\beta)$ لكل
$\tuple{\beta, \gamma} \in \cardfont{b} \times \cardfont{c}$.
الخريطة $f \mapsto f^*$ هي !!a{bijection} من
$\funfromto{\cardfont{c}}{(\funfromto{\cardfont{b}}{\cardfont{a}})}$
إلى $\funfromto{\cardfont{b} \times \cardfont{c}}{\cardfont{a}}$.
\end{proof}

ما \emph{نريده} هو طريقة سهلة لحساب
$\cardexpo{\cardfont{a}}{\cardfont{b}}$ حين نتعامل مع كاردينالات غير
متناهية. والنتيجة الآتية خطوة حسنة في هذا الاتجاه:

\begin{prop}\ollabel{cardexpo2reduct}
إذا كان $2 \leq \cardfont{a} \leq \cardfont{b}$ وكان
$\cardfont{b}$ غير متناهٍ، فإن
$\cardexpo{\cardfont{a}}{\cardfont{b}} =
\cardexpo{2}{\cardfont{b}}$.
\end{prop}

\begin{proof}
\begin{align*}
	\cardexpo{2}{\cardfont{b}} &\leq
	\cardexpo{\cardfont{a}}{\cardfont{b}}\text{، لأن $2 \leq \cardfont{a}$}\\
	&\leq \cardexpo{(\cardexpo{2}{\cardfont{a}})}{\cardfont{b}}
	\text{، بحسب \olref[opps]{cantorcor}}\\
	&= \cardexpo{2}{\cardfont{a} \cardtimes \cardfont{b}}
	\text{، بحسب \olref{simplecardexpo}} \\
	&= \cardexpo{2}{\cardfont{b}}
	\text{، بحسب \olref[simp]{cardplustimesmax}}.
\end{align*}
\end{proof}

لا ينبغي أن نتوقع حقاً أن نبسط هذا أكثر، لأن
$\cardfont{b} < \cardexpo{2}{\cardfont{b}}$ بحسب
\olref[card-arithmetic][opps]{cantorcor}. لكن هذا لا يخبرنا بما ينبغي
قوله عن $\cardexpo{\cardfont{a}}{\cardfont{b}}$ عندما
$\cardfont{b} < \cardfont{a}$. وبالطبع، إذا كان $\cardfont{b}$
\emph{متناهياً} فنحن نعرف ما نفعل.

\begin{prop}
إذا كان $\cardfont{a}$ غير متناهٍ وكان $0<n<\omega$، فإن
$\cardexpo{\cardfont{a}}{n} = \cardfont{a}$.
\end{prop}

\begin{proof}
$\cardexpo{\cardfont{a}}{n} = \cardfont{a} \cardtimes \cardfont{a}
\cardtimes \ldots \cardtimes \cardfont{a} = \cardfont{a}$، بحسب
\olref[simp]{cardplustimesmax}.
\end{proof}
\noindent
وفوق ذلك يمكننا، في بعض الحالات الأخرى، ضبط حجم
$\cardexpo{\cardfont{a}}{\cardfont{b}}$:

\begin{prop}
إذا كان $2 \leq \cardfont{b} < \cardfont{a} \leq
\cardexpo{2}{\cardfont{b}}$ وكان $\cardfont{b}$ غير متناهٍ، فإن
$\cardexpo{\cardfont{a}}{\cardfont{b}} =
\cardexpo{2}{\cardfont{b}}$.
\end{prop}

\begin{proof}
$\cardexpo{2}{\cardfont{b}}\leq
\cardexpo{\cardfont{a}}{\cardfont{b}} \leq
\cardexpo{(\cardexpo{2}{\cardfont{b}})}{\cardfont{b}} =
\cardexpo{2}{\cardfont{b}\cardtimes\cardfont{b}} =
\cardexpo{2}{\cardfont{b}}$، باستدلال مماثل لما في
\olref{cardexpo2reduct}.
\end{proof}
\noindent
أما بعد هذه النقطة، فتصير الأمور أشد دقة بكثير.

\end{document}
